
为陈述本小节的主要结果，
还需一些记号
（亦见第~\ref{subsec:3-8-4} 小节）。
$C^{\mathrm{amb}}$ 上的上边界算子
$\delta$ 在
$\widehat E(C^{\mathrm{amb}}[2])$
上诱导余导子 $\widehat\delta^{\circ}$：
\[
\begin{aligned}
\widehat\delta^{\circ}
(y_1\otimes\cdots\otimes y_k)
=\sum_{i=1}^{k}
(-1)^{\deg y_1+\cdots+\deg y_i}
y_1\otimes\cdots\otimes
\delta y_i\otimes\cdots\otimes y_k.
\end{aligned}
\]
它正是式~\eqref{eq:7-4-20}
中 $\mathfrak l_1$ 诱导的余导子。

令
$\mathbf y\in E(C^{\mathrm{amb}}[2])$、
$\mathbf x\in B(C[1])$，
记
\[
\begin{aligned}
\Delta^{k-1}\mathbf y
&=\sum_{c_1\in C_1}
\mathbf y_{c_1}^{(k;1)}
\otimes\cdots\otimes
\mathbf y_{c_1}^{(k;k)},\\
\Delta^{k-1}\mathbf x
&=\sum_{c_2\in C_2}
\mathbf x_{c_2}^{(k;1)}
\otimes\cdots\otimes
\mathbf x_{c_2}^{(k;k)}.
\end{aligned}
\]
这里 $C_1$ 依赖于 $k,\mathbf y$，
$C_2$ 依赖于 $k,\mathbf x$。
设 $C'_1\subset C_1$ 由满足
\[
\mathbf y_{c_1}^{(k;j)}
\notin E_0(C^{\mathrm{amb}}[2]),
\qquad j=1,\ldots,k
\]
的指标 $c_1$ 组成。

\begin{proposition}{7.4.34}\label{prop:7-4-34}
$\widehat{\mathfrak q}^{\circ}$
是严格滤过 $L_\infty$ 同态，
当且仅当 $\mathfrak q^{\circ}$
满足等式
\begin{equation}
\begin{aligned}
&(-1)^{\deg\mathbf y}
\mathfrak q^{\circ}_{\ell,k}
\bigl(\widehat\delta^{\circ}
(\mathbf y)\otimes\mathbf x\bigr)\\
&\quad+\sum_{c_1\in C_1,\,c_2\in C_2}
(-1)^{\substack{
\deg\mathbf y_{c_1}^{(2;2)}
\deg'\mathbf x_{c_2}^{(3;1)}
+\deg'\mathbf x_{c_2}^{(3;1)}\\
+\deg\mathbf y_{c_1}^{(2;1)}
+\deg\mathbf y_{c_1}^{(2;1)}
\deg\mathbf y_{c_1}^{(2;2)}}}\\
&\qquad
\mathfrak q^{\circ}\Bigl(
\mathbf y_{c_1}^{(2;1)}
\otimes\bigl(
\mathbf x_{c_2}^{(3;1)}\otimes\\
&\hspace{6em}
\mathfrak q^{\circ}
(\mathbf y_{c_1}^{(2;2)}
\otimes\mathbf x_{c_2}^{(3;2)})
\otimes\mathbf x_{c_2}^{(3;3)}
\bigr)\Bigr)=0.
\end{aligned}
\tag{7.4.35}\label{eq:7-4-35}
\end{equation}
在式~\eqref{eq:7-4-35}
及下文中，
我们省去诸如
$\mathfrak q^{\circ}_{
\ell(c_1),k(c_2)}
(\mathbf y_{c_1}^{(2;2)}
\otimes\mathbf x_{c_2}^{(3;2)})$
中的下标 $\ell(c_1),k(c_2)$\translatornote{源书把
$\mathbf x_{c_2}^{(3;2)}$
写成属于
$E_{k(c_2)}(C^{\mathrm{amb}}[2])$；
但 $\mathbf x$ 属于 $B(C[1])$，
式 (7.4.32) 的第二输入槽
亦为 $B_k(C[1])$。
这里据此更正，见疑误 A-0564。}；
这里
\[
\mathbf y_{c_1}^{(2;2)}
\in E_{\ell(c_1)}
(C^{\mathrm{amb}}[2]),
\qquad
\mathbf x_{c_2}^{(3;2)}
\in B_{k(c_2)}(C[1]).
\]
\end{proposition}

\begin{proof}
$\widehat{\mathfrak q}^{\circ}$
是滤过 $L_\infty$ 同态的条件
可写为
\begin{equation}
\begin{aligned}
\widehat{\mathfrak q}^{\circ}
(\mathfrak l_1(\mathbf y))
={}&\sum_{c_1\in C'_1}
\mathfrak l_2\bigl(
\widehat{\mathfrak q}^{\circ}
(\mathbf y_{c_1}^{(2;1)}),
\widehat{\mathfrak q}^{\circ}
(\mathbf y_{c_1}^{(2;2)})\bigr)\\
&+\mathfrak l_1\bigl(
\widehat{\mathfrak q}^{\circ}
(\mathbf y)\bigr),
\end{aligned}
\tag{7.4.36}\label{eq:7-4-36}
\end{equation}
因为当 $k\ne1$ 且
$\mathbf y\in E_k(C^{\mathrm{amb}}[2])$
时有 $\mathfrak l_k(\mathbf y)=0$。
注意式~\eqref{eq:7-4-36}
中的求和范围是 $c_1\in C'_1$。

由式~\eqref{eq:7-4-28-prime}、
后文的例 A3.4
及 $\Delta$ 的分次余交换性
\begin{equation}
\sum_{c_1}
\mathbf y_{c_1}^{(2;1)}
\otimes\mathbf y_{c_1}^{(2;2)}
=\sum_{c_1}
(-1)^{
\deg\mathbf y_{c_1}^{(2;1)}
\deg\mathbf y_{c_1}^{(2;2)}}
\mathbf y_{c_1}^{(2;2)}
\otimes\mathbf y_{c_1}^{(2;1)},
\tag{7.4.37}\label{eq:7-4-37}
\end{equation}
可得下面的等式。
\begin{equation}
\begin{aligned}
&\sum_{c_1\in C'_1}
\mathfrak l_2\bigl(
\widehat{\mathfrak q}^{\circ}
(\mathbf y_{c_1}^{(2;1)}),
\widehat{\mathfrak q}^{\circ}
(\mathbf y_{c_1}^{(2;2)})\bigr)
(\mathbf x)\\
&=\frac12
\sum_{\substack{c_1\in C'_1\\c_2\in C_2}}
(-1)^{\ast_1(i)}
\widehat{\mathfrak q}^{\circ}
\Bigl(\mathbf y_{c_1}^{(2;1)}
\otimes\bigl(
\mathbf x_{c_2}^{(3;1)}\otimes\\
&\hspace{7em}
\widehat{\mathfrak q}^{\circ}
(\mathbf y_{c_1}^{(2;2)}
\otimes\mathbf x_{c_2}^{(3;2)})
\otimes\mathbf x_{c_2}^{(3;3)}
\bigr)\Bigr)\\
&\quad+\frac12
\sum_{\substack{c_1\in C'_1\\c_2\in C_2}}
(-1)^{\ast_2(i)}
\widehat{\mathfrak q}^{\circ}
\Bigl(\mathbf y_{c_1}^{(2;2)}
\otimes\bigl(
\mathbf x_{c_2}^{(3;1)}\otimes\\
&\hspace{7em}
\widehat{\mathfrak q}^{\circ}
(\mathbf y_{c_1}^{(2;1)}
\otimes\mathbf x_{c_2}^{(3;2)})
\otimes\mathbf x_{c_2}^{(3;3)}
\bigr)\Bigr)\\
&=\frac12
\sum_{\substack{c_1\in C'_1\\c_2\in C_2}}
(-1)^{\ast_1(i)}
\widehat{\mathfrak q}^{\circ}
\Bigl(\mathbf y_{c_1}^{(2;1)}
\otimes\bigl(
\mathbf x_{c_2}^{(3;1)}\otimes\\
&\hspace{7em}
\widehat{\mathfrak q}^{\circ}
(\mathbf y_{c_1}^{(2;2)}
\otimes\mathbf x_{c_2}^{(3;2)})
\otimes\mathbf x_{c_2}^{(3;3)}
\bigr)\Bigr)\\
&\quad+\frac12
\sum_{\substack{c_1\in C'_1\\c_2\in C_2}}
(-1)^{\ast_3(i)}
\widehat{\mathfrak q}^{\circ}
\Bigl(\mathbf y_{c_1}^{(2;1)}
\otimes\bigl(
\mathbf x_{c_2}^{(3;1)}\otimes\\
&\hspace{7em}
\widehat{\mathfrak q}^{\circ}
(\mathbf y_{c_1}^{(2;2)}
\otimes\mathbf x_{c_2}^{(3;2)})
\otimes\mathbf x_{c_2}^{(3;3)}
\bigr)\Bigr)\\
&=\sum_{\substack{c_1\in C'_1\\c_2\in C_2}}
(-1)^{\ast_1(i)}
\widehat{\mathfrak q}^{\circ}
\Bigl(\mathbf y_{c_1}^{(2;1)}
\otimes\bigl(
\mathbf x_{c_2}^{(3;1)}\otimes\\
&\hspace{7em}
\widehat{\mathfrak q}^{\circ}
(\mathbf y_{c_1}^{(2;2)}
\otimes\mathbf x_{c_2}^{(3;2)})
\otimes\mathbf x_{c_2}^{(3;3)}
\bigr)\Bigr).
\end{aligned}
\tag{7.4.38.1}\label{eq:7-4-38-1}
\end{equation}
说明上述计算中的符号。
注意
$\deg\widehat{\mathfrak q}^{\circ}
(\mathbf y)
\equiv\deg(\mathbf y)+1\pmod2$。
第一个符号为
\[
\begin{aligned}
\ast_1(i)
={}&\deg\mathbf y_{c_1}^{(2;2)}
\deg'\mathbf x_{c_2}^{(3;1)}
+\deg'\mathbf x_{c_2}^{(3;1)}\\
&+(\deg\mathbf y_{c_1}^{(2;1)}+1)
\deg\mathbf y_{c_1}^{(2;2)}+1.
\end{aligned}
\]
末尾的 $+1$
来自式~\eqref{eq:7-4-28} 中的负号。
第二个符号为
\[
\begin{aligned}
\ast_2(i)
={}&\deg\mathbf y_{c_1}^{(2;1)}
\deg'\mathbf x_{c_2}^{(3;1)}
+\deg'\mathbf x_{c_2}^{(3;1)}\\
&+(\deg\mathbf y_{c_1}^{(2;1)}+1)
\deg\mathbf y_{c_1}^{(2;2)}\\
&+(\deg\mathbf y_{c_1}^{(2;1)}+1)
(\deg\mathbf y_{c_1}^{(2;2)}+1).
\end{aligned}
\]
令 $\ast'_2(i)$ 为把上式中的
$\mathbf y_{c_1}^{(2;1)}$
与 $\mathbf y_{c_1}^{(2;2)}$
互换后得到的式子。
第二个等号使用
式~\eqref{eq:7-4-37}；
此时
\[
\ast_3(i)
=\ast'_2(i)
+\deg\mathbf y_{c_1}^{(2;1)}
\deg\mathbf y_{c_1}^{(2;2)}
=\ast_1(i).
\]
于是第三个等号成立。

另一方面，由式~\eqref{eq:7-4-27-prime}
与例 A3.4 得
\[
\begin{aligned}
&(-1)^{\deg\mathbf y+1}
\mathfrak l_1\bigl(
\widehat{\mathfrak q}^{\circ}
(\mathbf y)\bigr)(\mathbf x)\\
&=\sum_{c_2\in C_2}
(-1)^{\deg\mathbf y\,\deg'\mathbf x_{c_2}^{(3;1)}
+\deg'\mathbf x_{c_2}^{(3;1)}}
\mathfrak m\Bigl(
\mathbf x_{c_2}^{(3;1)}
\otimes
\widehat{\mathfrak q}^{\circ}
(\mathbf y\otimes\mathbf x_{c_2}^{(3;2)})
\otimes\mathbf x_{c_2}^{(3;3)}
\Bigr)\\
&\quad+
\sum_{c_2\in C_2}
(-1)^{\deg\mathbf y+
\deg'\mathbf x_{c_2}^{(3;1)}}
\widehat{\mathfrak q}^{\circ}
\Bigl(\mathbf y\otimes\bigl(
\mathbf x_{c_2}^{(3;1)}
\otimes\mathfrak m
(\mathbf x_{c_2}^{(3;2)})
\otimes\mathbf x_{c_2}^{(3;3)}
\bigr)\Bigr).
\end{aligned}
\]
这里再次使用
$\deg\widehat{\mathfrak q}^{\circ}
(\mathbf y)
\equiv\deg(\mathbf y)+1\pmod2$。
因此，由式~\eqref{eq:7-4-33}
以及恒等式
\[
\sum_{c_1\in C_1\setminus C'_1}
\mathbf y_{c_1}^{(2;1)}
\otimes\mathbf y_{c_1}^{(2;2)}
=1\otimes\mathbf y
+\mathbf y\otimes1
\]
得到
\begin{equation}
\begin{aligned}
&(-1)^{\deg\mathbf y+1}
\mathfrak l_1\bigl(
\widehat{\mathfrak q}^{\circ}
(\mathbf y)\bigr)(\mathbf x)\\
&=\sum_{\substack{
c_1\in C_1\setminus C'_1\\
c_2\in C_2}}
(-1)^{\substack{
\deg\mathbf y_{c_1}^{(2;2)}
\deg'\mathbf x_{c_2}^{(3;1)}
+\deg'\mathbf x_{c_2}^{(3;1)}\\
+\deg\mathbf y_{c_1}^{(2;1)}
(\deg\mathbf y_{c_1}^{(2;2)}+1)}}\\
&\qquad
\widehat{\mathfrak q}^{\circ}
\Bigl(\mathbf y_{c_1}^{(2;1)}
\otimes\bigl(
\mathbf x_{c_2}^{(3;1)}\otimes\\
&\hspace{6em}
\widehat{\mathfrak q}^{\circ}
(\mathbf y_{c_1}^{(2;2)}
\otimes\mathbf x_{c_2}^{(3;2)})
\otimes\mathbf x_{c_2}^{(3;3)}
\bigr)\Bigr).
\end{aligned}
\tag{7.4.38.2}\label{eq:7-4-38-2}
\end{equation}
最后，
式~\eqref{eq:7-4-35}
由式~\eqref{eq:7-4-36}、
\eqref{eq:7-4-38-1}、
\eqref{eq:7-4-38-2}
以及模 $2$ 的同余
\[
(\deg\mathbf y_{c_1}^{(2;1)}+1)
\deg\mathbf y_{c_1}^{(2;2)}
+\deg\mathbf y
\equiv
\deg\mathbf y_{c_1}^{(2;1)}
+\deg\mathbf y_{c_1}^{(2;1)}
\deg\mathbf y_{c_1}^{(2;2)}
\pmod2
\]
推出。
\end{proof}
